\documentclass[10pt,letterpaper]{article}
\usepackage[T1]{fontenc}
\usepackage[utf8]{inputenc}
\usepackage[spanish,es-tabla]{babel}
\usepackage{csquotes}
\usepackage{mathpazo}
\usepackage[scaled=0.9]{helvet}
\usepackage{courier}
\usepackage{calc}
\usepackage[top=2cm, bottom=2cm, left=2cm, right=2cm]{geometry}
\usepackage{xcolor}
\usepackage{booktabs}
\usepackage{longtable}
\usepackage{array}
\usepackage{ragged2e}
\usepackage{xurl}
\usepackage[strings]{underscore}
\usepackage{fancyhdr}
\usepackage{sectsty}
\usepackage{tocloft}
\usepackage[colorlinks=true,linkcolor=blue,urlcolor=blue,citecolor=blue,breaklinks=true]{hyperref}
\usepackage{enumitem}

\definecolor{uncpblue}{HTML}{1A237E}
\allsectionsfont{\color{uncpblue}}

\renewcommand{\cftsecleader}{\cftdotfill{\cftsecdotsep}}
\cftpagenumbersoff{section}
\cftpagenumbersoff{subsection}
\cftpagenumbersoff{subsubsection}
\renewcommand{\cftsecfont}{\normalfont}
\renewcommand{\cftsubsecfont}{\normalfont}
\renewcommand{\cftsubsubsecfont}{\normalfont}
\setlength{\parindent}{0pt}
\setlength{\parskip}{4pt}
\setlength{\tabcolsep}{3pt}
\renewcommand{\arraystretch}{0.85}

\pagestyle{fancy}
\fancyhf{}
\fancyhead[L]{\small\textcolor{uncpblue}{\textbf{SGSI-UNCP}}}
\fancyhead[R]{\small\textcolor{uncpblue}{D-SGSI-03}}
\fancyfoot[C]{\thepage}
\renewcommand{\headrulewidth}{0.4pt}
\renewcommand{\footrulewidth}{0.4pt}
\setlength{\headheight}{14pt}

\setlist{nosep,leftmargin=*,after=\vspace{2pt}}
\setlength{\emergencystretch}{2em}

\newcommand{\makecover}{
\begin{titlepage}
\centering
\vspace*{2cm}
{\Huge\bfseries\color{uncpblue} SGSI-UNCP\par}
\vspace{0.7cm}
{\LARGE Política General de Seguridad de la Información\par}
\vspace{0.5cm}
{\large Documento Base del Sistema de Gestión\\de Seguridad de la Información\par}
\vspace{1.5cm}
{\small Universidad Nacional del Centro del Perú\par}
\vfill
\end{titlepage}
}

\begin{document}

\makecover
\setcounter{tocdepth}{2}
\RaggedRight
\tableofcontents
\newpage

\RaggedRight
\section{Política General de Seguridad de la
Información}\label{poluxedtica-general-de-seguridad-de-la-informaciuxf3n}

\textbf{Organización:} Universidad Nacional del Centro del Perú (UNCP)
\textbf{Aprobado por:} Rector / Presidente del Comité de Gobierno
Digital \textbf{Norma:} ISO/IEC 27001:2022 (Cláusula 5.2)

\begin{center}\rule{0.5\linewidth}{0.5pt}\end{center}

\subsection{Declaración de
Compromiso}\label{declaraciuxf3n-de-compromiso}

La Universidad Nacional del Centro del Perú (UNCP), consciente de que la
información es un activo crítico para el cumplimiento de sus fines
misionales ---formación académica, investigación científica y
responsabilidad social--- y en estricta alineación con el Plan de
Gobierno y Transformación Digital (PGTD 2026-2030), la Política Nacional
de Transformación Digital al 2030 (PNTD) y la Política General de
Gobierno 2025-2026 (D.S. N.° 141-2025-PCM), establece la presente
Política General de Seguridad de la Información.

La Alta Dirección, representada por el Rectorado y el Comité de Gobierno
Digital (CGD), se compromete a liderar, respaldar y dotar de los
recursos necesarios para la implementación, mantenimiento y mejora
continua del Sistema de Gestión de Seguridad de la Información (SGSI),
garantizando la confidencialidad, integridad y disponibilidad de la
información institucional, protegiendo así los intereses de estudiantes
(≈11,700), docentes, personal administrativo y la sociedad en general.

\subsection{Objetivos de Seguridad de la
Información}\label{objetivos-de-seguridad-de-la-informaciuxf3n}

El SGSI de la UNCP se alinea al Objetivo de Gobierno y Transformación
Digital 3 (OGTD3 --- Seguridad y Confianza Digital) y persigue los
siguientes objetivos específicos:

\begin{enumerate}
\def\labelenumi{\arabic{enumi}.}
\item
  \textbf{Confidencialidad:} Garantizar que la información académica,
  personal y médica (datos sensibles según Ley N.° 29733) sea accesible
  únicamente al personal autorizado, mediante controles de acceso
  basados en roles y cifrado de datos en tránsito y reposo.
\item
  \textbf{Integridad:} Proteger la exactitud y completitud de los
  registros académicos (notas, grados, títulos), información financiera
  y expedientes institucionales contra modificaciones no autorizadas,
  asegurando la confiabilidad de los procesos digitales mediante firmas
  electrónicas y controles de cambio.
\item
  \textbf{Disponibilidad:} Asegurar que los servicios tecnológicos
  críticos ---Campus Virtual (Moodle 4.1), ERP ADESA, correo
  institucional (Microsoft 365) y portal web--- operen con los niveles
  de disponibilidad establecidos en el PGTD (mínimo 99.5\% en horario
  académico), respaldando la continuidad operativa de la universidad.
\item
  \textbf{Cumplimiento Normativo:} Cumplir estrictamente con la Ley N.°
  29733 (Protección de Datos Personales), el D.L. N.° 1412 (Gobierno
  Digital), el D.S. N.° 029-2021-PCM (Marco de Confianza Digital), la
  Ley de Ciberdefensa, y los requisitos de SUNEDU aplicables a los
  sistemas de información universitarios.
\item
  \textbf{Cultura de Seguridad:} Desarrollar un nivel óptimo de cultura
  digital (OGTD5 y OGTD6), concienciando a toda la comunidad
  universitaria sobre sus responsabilidades en ciberseguridad mediante
  programas anuales de sensibilización y capacitación obligatoria.
\end{enumerate}

\subsection{Alcance}\label{alcance}

Esta política es de cumplimiento obligatorio para:

\begin{itemize}
\tightlist
\item
  Todo el personal docente, administrativo, autoridades y estudiantes de
  la UNCP, incluyendo a aquellos que accedan a los recursos informáticos
  institucionales de forma remota.
\item
  Proveedores de servicios tecnológicos (ej. Huawei Cloud, Microsoft
  365), contratistas y terceros con acceso a los activos de información
  de la UNCP.
\item
  Todos los procesos, servicios e infraestructuras contemplados en el
  documento D-SGSI-02 (Alcance del SGSI), abarcando las cuatro sedes
  universitarias y los puntos de acceso remoto autorizados.
\end{itemize}

\subsection{Principios Rectores}\label{principios-rectores}

\begin{enumerate}
\def\labelenumi{\arabic{enumi}.}
\item
  \textbf{Gestión de Riesgos:} Las medidas de seguridad se adoptarán con
  base en la evaluación periódica de riesgos (ISO 31000), buscando
  mitigarlos hasta un nivel aceptable para la UNCP. La metodología de
  evaluación de riesgos se describe en el documento D-SGSI-04.
\item
  \textbf{Seguridad desde el Diseño:} Todo nuevo proyecto tecnológico
  definido en el PGTD debe incorporar requisitos de seguridad y
  privacidad desde su fase de concepción, conforme al principio de
  \emph{Privacy by Design} establecido en la Ley N.° 29733.
\item
  \textbf{Gestión de Incidentes:} Todo usuario tiene la obligación de
  reportar de inmediato cualquier evento, debilidad o incidente de
  seguridad sospechoso a través del canal oficial establecido (CSIRT
  UNCP / Mesa de Ayuda OTI), conforme al procedimiento P-SGSI-07.
\item
  \textbf{Sanciones:} El incumplimiento de esta política y de los
  controles del SGSI estará sujeto a medidas disciplinarias conforme al
  Reglamento Interno de la UNCP y la legislación vigente, sin perjuicio
  de las responsabilidades civiles o penales aplicables.
\item
  \textbf{Mejora Continua:} El SGSI opera bajo el ciclo PHVA
  (Planificar-Hacer-Verificar-Actuar), realizando auditorías internas
  anuales y revisiones por la dirección para asegurar su eficacia
  continua.
\end{enumerate}

\subsection{Clasificación de la
Información}\label{clasificaciuxf3n-de-la-informaciuxf3n}

La información institucional se clasifica en cuatro niveles, conforme al
procedimiento P-SGSI-05:

{\setlength{\RaggedRightRightskip}{0pt plus 5em}\small\sloppy \begin{longtable}[]{@{}
  >{\hspace{0pt}\RaggedRight\arraybackslash}p{(\linewidth - 6\tabcolsep) * \real{0.1458}}
  >{\hspace{0pt}\RaggedRight\arraybackslash}p{(\linewidth - 6\tabcolsep) * \real{0.2292}}
  >{\hspace{0pt}\RaggedRight\arraybackslash}p{(\linewidth - 6\tabcolsep) * \real{0.2083}}
  >{\hspace{0pt}\RaggedRight\arraybackslash}p{(\linewidth - 6\tabcolsep) * \real{0.4167}}@{}}
\toprule\noalign{}
\begin{minipage}[b]{\linewidth}\raggedright
Nivel
\end{minipage} & \begin{minipage}[b]{\linewidth}\raggedright
Categoría
\end{minipage} & \begin{minipage}[b]{\linewidth}\raggedright
Ejemplos
\end{minipage} & \begin{minipage}[b]{\linewidth}\raggedright
Controles Mínimos
\end{minipage} \\
\midrule\noalign{}
\endhead
\bottomrule\noalign{}
\endlastfoot
4 & Confidencial & Datos personales sensibles (salud, disciplina),
secretos industriales, estrategias institucionales & Cifrado AES-256,
acceso restringido por rol, auditoría de accesos \\
3 & Interna Restringida & Expedientes académicos, contratos, planillas,
resoluciones rectorales & Control de acceso, logs de auditoría, copias
de seguridad cifradas \\
2 & Interna & Comunicados, directivas, manuales, material didáctico &
Acceso autenticado, prevención de fuga de datos \\
1 & Pública & Información institucional publicable, portal web, datos
abiertos (PNDA) & Integridad de publicación, disponibilidad del
portal \\
\end{longtable}\normalsize}

\subsection{Responsabilidades}\label{responsabilidades}

\begin{itemize}
\tightlist
\item
  \textbf{Comité de Gobierno Digital (CGD):} Aprobar y revisar
  anualmente esta política, asegurando su alineación con los objetivos
  estratégicos de la UNCP y el PGTD.
\item
  \textbf{Oficial de Seguridad y Confianza Digital:} Diseñar,
  implementar, operar y evaluar el SGSI, reportando su desempeño a la
  Alta Dirección y coordinando con la OTI la ejecución de controles
  técnicos.
\item
  \textbf{Oficina de Tecnologías de la Información (OTI):} Ejecutar los
  controles técnicos y operativos definidos por el SGSI, administrar la
  infraestructura tecnológica y atender incidentes de seguridad.
\item
  \textbf{Dueños de Procesos / Activos:} Identificar, clasificar y
  proteger la información bajo su responsabilidad, participando
  activamente en las evaluaciones de riesgo de sus procesos.
\item
  \textbf{Usuarios (Comunidad Universitaria):} Leer, comprender y
  aplicar las normativas de seguridad en su trabajo diario, reportando
  cualquier anomalía o incidente de seguridad detectado.
\end{itemize}

\subsection{Revisión y Mejora
Continua}\label{revisiuxf3n-y-mejora-continua}

La presente política será revisada al menos una vez al año por el Comité
de Gobierno Digital o cuando ocurran cambios significativos en el
entorno tecnológico, normativo o estratégico de la UNCP ---tales como
modificaciones a la Ley N.° 29733, entrada en vigor de nuevas
disposiciones del Marco de Confianza Digital, o cambios mayores en la
infraestructura tecnológica--- con el fin de asegurar su conveniencia,
adecuación y eficacia continua.

Los resultados de la revisión se documentarán en el acta del CGD y, de
ser necesario, se generará una nueva versión de esta política siguiendo
el procedimiento de control documental establecido en R-SGSI-00.

\begin{center}\rule{0.5\linewidth}{0.5pt}\end{center}

\textbf{Firmas de Aprobación:}

\begin{center}\rule{0.5\linewidth}{0.5pt}\end{center}

\textbf{Rector de la UNCP} Presidente del Comité de Gobierno Digital

\begin{center}\rule{0.5\linewidth}{0.5pt}\end{center}

\textbf{Director(a) General de Administración} Líder de Gobierno y
Transformación Digital

\begin{center}\rule{0.5\linewidth}{0.5pt}\end{center}

\textbf{Oficial de Seguridad y Confianza Digital}

\end{document}
